%%%%%
%
%% Dans ce fichier, les �ventuels accents sont cod�s
%% en ISO-8859-15 (aka latin9)
%% La ligne suivante donne trois exemples 
%% e accent aigu : �, e accent grave : �, e-dans-l'o : �
\section{Presence d'un top de hauteur de base de nuage < a une valeur de seuil}

\itemstyle{Op\'erateur~:} \operfmt{PresenceNuageEch\_tps\_bnds}
\index{PresenceNuageEch\_tps\_bnds}

%%%%%OBJET%%%%
\itemstyle{Objet~:}
Cette fonction calcule, pour chaque instant $t$ d'une
base de temps $time_{out}$, le pourcentage de top de hauteur de base de nuage 
 inf�rieur � une valeur de seuil pour au moins une des 3 couches de hauteur de base de nuage $X_{i_{in}}$ 
entre les bornes donn�es par le $time_{bnds}$ associ� � $time_{out}$.
Aucune supposition n'est faite sur la r�gularit� de
l'�chantillonnage des 2 variables. Il est cependant n�cessaire
que $time_{in}$ (le temps des $X_{i_{in}}$) soit "� peu pr�s" tri� par ordre
croissant.


%%%%%ARGUMENTS%%%%
\bigskip
%% Dans ce fichier, les �ventuels accents sont cod�s
%% en ISO-8859-15 (aka latin9)
%% La ligne suivante donne trois exemples 
%% e accent aigu : �, e accent grave : �, e-dans-l'o : �
%%%%%ENTREES%%%%
\begin{tabarg}
\multicolumn{3}{l}{\titretab {Entr\'ees}}\\
\hline
$X1_{in}$ & Vecteur, \texttt{float} & Hauteur de la base des nuages de la premi�re couche\\
$X2_{in}$ & Vecteur, \texttt{float} & Hauteur de la base des nuages de la deuxi�me couche\\
$X3_{in}$ & Vecteur, \texttt{float} & Hauteur de la base des nuages de la troisi�me couche\\
$toler_{mnq}$ & Coeff., \texttt{Int\_32} & Nombre de top minimum pour faire le calcul (nb. �chant). A chaque instant de la base de temps d'arriv�e, 
le nombre de top est calcul� si le nombre d'�chantillons dans l'intervalle de temps est strictement sup�rieur � $toler_{mnq}$. Sinon, 
le r�sultat est "valeur manquante". Si $toler_{mnq}$ = 0, le calcul est fait d�s qu'il y a au moins une valeur valide dans l'�chantillon. \\
$time_{out}$ & Vecteur, \texttt{double} & Base de temps d'arriv�e (secondes)\\
$val_{inv}$ & Coeff., \texttt{float} & valeur invalide pour laquelle le nombre de top n'est pas calcul�. \\
\end{tabarg}

\medskip
%%%%%SORTIES%%%%
\begin{tabarg}
\multicolumn{3}{l}{\titretab {Sorties}}\\
\hline
$X_{out}$ & Vecteur, \texttt{float} & $X_{in}$ Pourcentage de valeurs pr�sentant au moins un top < $val_{inv}$ 
pour une des 3 hauteurs de hbn mesur�e\\
\end{tabarg}



%%%%%FORMULE%%%%
\itemstyle{Formule}

%% Formule de l'operateur : PresenceNuageEch_tps_bnds
%% Dans ce fichier, les �ventuels accents sont cod�s
%% en ISO-8859-15 (aka latin9)
%% La ligne suivante donne trois exemples 
%% e accent aigu : �, e accent grave : �, e-dans-l'o : �

Pour chaque couche Xn:

Pour  i=$0$, Nbre  $Time_{bnds}$ :

%Pour chaque j entre $t_b[i][0]$ et $t_b[i][1]$:

\[
X_{out}[i] =
\frac{\displaystyle{\sum_{tb[i][0]}^{tb[i][1]} (top < val_{inv} ) * 100} }{N_i}
\]
%\frac{\displaystyle{\sum_{tb[i][0]}^{tb[i][1]} Xn_{in}[j]} }{Nn_i}$
%$\frac{\displaystyle{\texttt{nb de }X_{in}[j]}}{N_i}$
%\frac{\displaystyle{\sum_{tb[i][0]}^{tb[i][1]} {\max_{{1}\lt i \lt {3} }(Xn_{in}[j] < val_{seuil}) }} }{Nn_i}
%\frac{\displaystyle{\sum_{tb[i][0]}^{tb[i][1]} \max_{{1}\lt i \lt {3} } \texttt{Xn\_{in}[j] < val\_{seuil}} } }{Nn_i}
\[
\texttt{Avec top = }
\begin{cases}
\begin{tabular}{ll}
\texttt{1}
&
\texttt{ si  au moins 1 $Xn_{in}[i]$ est inf\'erieure \`a $val_{inv}$ } \\
\texttt{0} & \texttt{ sinon} \\
\end{tabular} 
\end{cases}
\]
Avec $N_i$ le nombre de donn\'ees dans l'intervalle tb .


%%%%COMMENTAIRE%%%%
\itemstyle{Commentaire~:}
\textit{tb} est un tableau d'intervalles de temps d�termin�s � partir du vecteur temps $time_{out}$.

%%%%%AUTEUR%%%%
\itemstyle{Auteur~:}
CNRM/GMEI/TRAMM

%%%%%REFERENCE%%%%
\itemstyle{R\'ef\'erences~:}
N�ant.
